\subsection{RQ1.2: Fine-Grained Refinement}
\label{sec:analysis-rq12}

RQ1.2 asks whether refining the carried-forward late-stage region from RQ1.1
with a block-level boundary changes the picture of which split is most
suitable under controlled local execution. The analytical contribution of
this section is not to repeat the per-condition latencies tabulated in
\cref{tab:rq12-results}, but to extract the one inference that RQ1.1 could
not make: at constant activation transfer, compute placement around the
boundary is an independently observable cost dimension. The RQ1.2 result is
therefore a controlled refinement of RQ1.1's joint reading rather than a
replacement of it; it confirms the carried-forward region as a flat-bottomed
valley and provides the controlled-contrast leg of the two-dimensional
explanation that \cref{sec:analysis-rq13} then synthesises. The design
itself --- staged refinement inside a carried-forward region rather than a
fresh global search --- follows the partitioning-literature practice of
treating split-point selection as a hierarchical procedure over a small set
of architecturally meaningful boundaries~\cite{dnnsplit,%
fine_grained_partitioning,hierarchical_partitioning}.

\paragraph{The refined region is a band, not a peak.}
The three refined split conditions fall in a 0.235~ms window
(\(80.157\)~ms, \(80.348\)~ms, \(80.392\)~ms in \cref{tab:rq12-results}),
all between \(+3.34\%\) and \(+3.64\%\) over the monolithic baseline. The
within-run per-iteration \(95\%\) confidence intervals on the mean are
\([79.998,\ 80.316]\) for \texttt{split\_after\_layer3\_block0},
\([80.192,\ 80.504]\) for \texttt{split\_after\_layer2}, and
\([80.238,\ 80.545]\) for \texttt{split\_after\_layer3}
(\texttt{condition\_summaries.csv}). All three are disjoint from the
monolithic CI of \([77.338,\ 77.792]\), so the split-versus-monolithic
overhead is robust within-run; the pairwise CIs \emph{among} the three
splits overlap, so the ordering between them is not separated at this
measurement resolution. The carry-forward rule in
\cref{sec:methods-carryforward} nevertheless fires deterministically: the
near-best \(5\%\) window admits all three splits as eligible (threshold
\(84.16\)~ms in \texttt{carry\_forward.json}) and the rule promotes the
raw-fastest mean to main and the next-lowest to reference. The defensible
reading is therefore that RQ1.2 confirms the carried-forward region as a
genuinely flat band rather than that it discovered a meaningfully better
split inside that band.

\paragraph{Equal activation size with unequal cost isolates compute placement.}
The instructive comparison is
\texttt{split\_after\_layer3\_block0} against \texttt{split\_after\_layer3}.
Both transfer exactly \(200{,}704\)~B (\(196\)~KiB) of activation in FP32,
yet their boundary-crossing intervals differ by \(7.81\)~ms
(\(33.97\)~ms versus \(26.16\)~ms in \texttt{cross\_condition.csv}). The
predeclared boundary-crossing definition in \cref{sec:methods-contract} is
a composite that bundles remote service-B compute with serialisation,
gRPC transport, and deserialisation. The per-segment compute decomposition
in \texttt{cross\_condition.csv} shows service-B mean compute of
\(31.87\)~ms for \texttt{split\_after\_layer3\_block0} versus
\(24.06\)~ms for \texttt{split\_after\_layer3} --- a \(7.81\)~ms gap --- and a
non-compute residual of approximately \(2.1\)~ms for both. The
boundary-crossing difference is therefore the service-B compute difference
on the wire-cost-equal pair, not a transport or serialisation
difference. Because activation transfer is held constant by construction,
the proximate explanatory factor must be the per-block compute distribution
around the cut. This is the cleanest within-run evidence the thesis offers
for treating compute placement as an independent dimension alongside
activation-transfer size, and it is consistent with the prior observation
that block-level execution time can vary across candidate boundaries even
when their activation footprints are similar~\cite{pareto_split}. The
composite-measurement framing also matches the per-layer profiling logic on
which Neurosurgeon's original split-point cost
model rests~\cite{neurosurgeon}.

\paragraph{The refinement does not reshuffle the carry-forward recommendation in any practically meaningful sense.}
The RQ1.2 carry-forward rule promotes
\texttt{split\_after\_layer3\_block0} to main and demotes
\texttt{split\_after\_layer3} from RQ1.1's main to RQ1.2's rejected
position (\texttt{carry\_forward.json}). At first glance this is a change of
recommendation; closer inspection shows it is a tie-break inside an
operationally indistinguishable set. The mean separation between the
promoted and demoted conditions is \(0.235\)~ms (\(0.29\%\)), smaller than
the per-condition round range of \texttt{split\_after\_layer3\_block0}
itself (\(1.46\)~ms across the five measured rounds in
\texttt{round\_consistency.json}). The block-level winner is also the
noisiest of the three splits round-to-round (round CV \(0.0070\), against
\(0.0046\) and \(0.0033\) for the two coarse anchors). The
thesis-facing reading is therefore not that RQ1.2 supersedes RQ1.1's split
choice. It is that the late-stage carried-forward region is an operationally
flat band of practically equivalent boundaries, that the carry-forward rule
selects deterministically within that band, and that any of the three
eligible splits is a defensible carry-forward candidate at this measurement
resolution.

\paragraph{Validation.}
Three pieces of evidence support treating the RQ1.2 within-run findings as
credible for the measured environment. First, functional parity is intact:
both local and gRPC parity checks report
\texttt{max\_abs\_diff} \(=0.0\) at the predeclared tolerance of
\(1.0\times10^{-5}\) for \texttt{layer2}, \texttt{layer3.0}, and
\texttt{layer3} (\texttt{parity\_validation.json}), so the \(7.81\)~ms
boundary-crossing gap between the two \(196\)~KiB conditions is not
confounded with numerical drift across the partition. Second, round-level
consistency is in line with RQ1.1: per-condition round CVs lie between
\(0.0017\) (monolithic) and \(0.0070\)
(\texttt{split\_after\_layer3\_block0}) in
\texttt{round\_consistency.json}, supporting the reporting practice of
combining per-iteration and round-level statistics rather than mean
alone~\cite{benchmarking_sc15}. Third, host controls and the measurement
contract are identical to RQ1.1 (\cref{tab:rq12-config}), so the comparison
with RQ1.1 is internally fair on infrastructure. Two boundaries to this
validation must be stated explicitly. The 0.235~ms span between the three
split means is narrower than the largest per-condition round range
(\(1.46\)~ms for the selected main), so the ordering among the three splits
is not robust round-to-round; only the split-versus-monolithic separation
is. And, as in RQ1.1, the run is single-host, single-build, fixed-seed,
CPU-only batch-one localhost gRPC, so cross-host or cross-seed replication
is not in scope.

\paragraph{Evaluation.}
The evaluative reading of RQ1.2 has two facets. Against the monolithic
baseline, refinement does not change the qualitative finding of RQ1.1:
every refined split is measurably slower (between \(+2.59\)~ms and
\(+2.83\)~ms; \(+3.34\%\) to \(+3.64\%\)), with the three split CIs
disjoint from the monolithic CI. Within the refined set, the evaluative
question is not which split is fastest --- the three are not separated at
the within-run per-iteration CI level --- but what the contrast reveals.
The controlled equal-activation pair attributes a \(7.81\)~ms
boundary-crossing difference essentially in full to per-block compute
distribution, which is the single piece of evidence the staged design was
constructed to produce. This is a methodological gain for RQ1.3 even
though it is not a latency gain for any deployment, and it justifies
carrying both \texttt{split\_after\_layer3\_block0} (the rule-selected
main) and \texttt{split\_after\_layer2} (the reference) into the
subsequent Kubernetes and Azure stages rather than committing to a single
boundary at this resolution.

\paragraph{Answer to RQ1.2.}
RQ1.2 provides direct but scoped evidence that block-level refinement
inside the carried-forward late-stage region from RQ1.1 does not overturn
the coarse-level recommendation in any practically meaningful sense, and
that the most useful refined finding is structural rather than rank-order.
The three refined splits cluster in a \(0.235\)~ms band whose internal
ordering is not separated at the per-iteration \(95\%\) CI level and is
narrower than the round-to-round range of the rule-selected main, so the
carry-forward outcome --- \texttt{split\_after\_layer3\_block0} as main and
\texttt{split\_after\_layer2} as reference, with \texttt{split\_after\_layer3}
rejected --- should be read as a deterministic rule output over
operationally equivalent options rather than as a separation. The
analytically load-bearing finding is the controlled
\texttt{split\_after\_layer3\_block0}~vs.~\texttt{split\_after\_layer3}
contrast at constant \(196\)~KiB activation: holding transfer constant by
construction, the \(7.81\)~ms boundary-crossing gap is essentially the
service-B compute gap, with a non-compute residual that is the same for
both conditions. Under RQ1.2's conditions, compute placement is therefore
an independently observable cost dimension at the boundary, consistent with
prior partitioning evidence that block-level execution time can vary at
equal activation size~\cite{pareto_split}. This is the controlled inference
that RQ1.1 could not make on a joint reading and is the empirical basis on
which \cref{sec:analysis-rq13} builds its two-dimensional explanation.
